\subsection{Measurable equivalence relations}

Given a topological space X and an equivalence relation R in X, we shall say that R is Hausdorff if the quotient space X/R is Hausdorff.

Recall (GT, I, §8, No.~3, Prop. 8) that when R is an open equivalence relation, it comes to the same to say that the graph of R in X × X is closed.

Let p be a mapping of X into a Hausdorff topological space B, and let R be the equivalence relation \(p(x) = p(y)\) in X; if K is a compact subset of X such that the restriction of p to K is continuous, then the relation \(R_{K}\) induced by R on K is Hausdorff, because the quotient space \(K/R_{K}\) is homeomorphic to the space \(p(K)\) , which is compact (GT, I, §9, No.~4, Th. 2 and its Cor. 2). If T is a locally compact space, \(\mu\) is a positive measure on T, and p is a \(\mu\) -measurable mapping of T into a Hausdorff topological space B, one thus sees that there exists a \(\mu\) -dense set (Ch. IV, §5, No.~8) of compact subsets K of T for which the relation \(R_{K}\) is Hausdorff. We are thus led to make the following definition:

DEFINITION 3. --- Let T be a locally compact space, \(\mu\) a positive measure on T. An equivalence relation R in T is said to be \(\mu\) -measurable if there exists a \(\mu\) -dense set of compact subsets K of T for which the relation \(R_{K}\) is Hausdorff.

If R is Hausdorff then R is \(\mu\) -measurable, because if \(\varphi\) is the canonical mapping of T onto the Hausdorff topological space T/R, \(\varphi\) is continuous and R is equivalent to \(\varphi(x)=\varphi(y)\) . Similarly, if R is such that the saturation for R of every compact subset of T is closed (in particular, if R is a closed equivalence relation), then R is \(\mu\) -measurable, because for every compact subset K of T, the relation \(R_{K}\) is closed, hence Hausdorff (GT, I, §10, No.~4, Prop. 8).

Note that if R is \(\mu\) -measurable, then R is also measurable for every measure on T with base \(\mu\) .

PROPOSITION 2. --- Let T be a locally compact space countable at infinity, \(\mu\) a positive measure on T.

\begin{enumerate}
\def\labelenumi{\arabic{enumi})}
\item
  For every \(\mu\) -measurable equivalence relation R in T, there exist a locally compact space B and a \(\mu\) -measurable mapping p of T into B such that R is equivalent to the relation \(p(x) = p(y)\) .
\item
  If, moreover, T admits a countable base, one can suppose that B admits a countable base.
\end{enumerate}

Since T is countable at infinity, there exists an increasing sequence \((\mathrm{K}_{n})_{n\geqslant1}\) of compact subsets of T such that T is the union of the \(K_{n}\) and a \(\mu\) -negligible set N, and such that each of the relations \(R_{K_{n}}\) is Hausdorff. Let \(B_{n}\) be the quotient space \(K_{n}/R_{K_{n}}\) , which is compact, and let \(B_{n}^{\prime}\) be the compact space that is the topological sum of \(B_{n}\) and a point \(a_{n}\) . Let \(q_{n}\) be the canonical mapping of \(K_{n}\) onto \(B_{n}\) ; we extend \(q_{n}\) to a mapping \(p_{n}\) of T into \(B_{n}^{\prime}\) in the following manner: if \(x\in T\) is equivalent mod R to an element \(y\in K_{n}\) , set \(p_{n}(x)=q_{n}(y)\) ; in the contrary case, set \(p_{n}(x)=a_{n}\) . Let \(B^{\prime}\) be the product space \(\prod_{i=1}^{\infty}B_{n}^{\prime}\) , which is compact, and let \(p^{\prime}\) be the mapping \(x \mapsto (p_n(x))\) of T into B'. Let us show that \(p'\) is \(\mu\) -measurable: it suffices (Ch. IV, §5, No.~3, Th. 1) to prove that each of the mappings \(p_n\) is measurable, and for this it suffices that the restriction of \(p_n\) to each \(K_m\) be measurable. Now, this is obvious if \(m \leqslant n\) ; if, on the contrary, m \textgreater{} n, let \(K_{nm}\) be the saturation of \(K_n\) for \(R_{K_m}\) , which is a compact subset of \(K_m\) (GT, I, §9, No.~4, Th. 2); since \(p_n\) is constant on \(K_m - K_{nm}\) , it suffices to prove that the restriction of \(p_n\) to \(K_{nm}\) is continuous, which is obvious on account of the canonical isomorphism between \(K_{nm}/R_{K_{nm}}\) and \(K_n/R_{K_n}\) (GT, I, §9, No.~4, Cor. 4 of Th. 2).

Let A be the saturation of \(\bigcup_{n} \mathrm{K}_{n}\) for the relation R, and let \(\mathrm{N}' = \mathrm{T} - \mathrm{A} \subset \mathrm{N}\) . We shall see that the relation \(p'(x) = p'(y)\) is equivalent to the relation \(\langle \mathrm{R}\{x, y\}\) or \((x, y) \in \mathrm{N}' \times \mathrm{N}' \rangle\) . For, if \(\mathrm{R}\{x, y\}\) then \(p_n(x) = p_n(y)\) for all \(n\) , therefore \(p'(x) = p'(y)\) ; and if \(x \in \mathrm{N}'\) , \(y \in \mathrm{N}'\) then \(p_n(x) = p_n(y) = a_n\) for all \(n\) , therefore \(p'(x) = p'(y)\) . If on the other hand \(x\) and \(y\) are in A and are not equivalent mod R, then there exist an integer \(n\) , an element \(x' \in \mathrm{K}_n\) (resp. \(y' \in \mathrm{K}_n\) ) equivalent mod R to \(x\) (resp. \(y\) ) such that \(x'\) is not equivalent to \(y'\) mod \(\mathrm{R}_{\mathrm{K}_n}\) ; therefore \(p_n(x) \neq p_n(y)\) , consequently \(p'(x) \neq p'(y)\) . Finally, if \(x \in \mathrm{N}'\) and \(y \in \mathrm{A}\) , then \(p_n(y) \in \mathrm{B}_n\) for \(n\) sufficiently large and \(p_n(x) = a_n\) for all \(n\) , therefore \(p'(x) \neq p'(y)\) , which establishes our assertion.

Consider then the quotient set \(B_{0} = N^{\prime}/R_{N^{\prime}}\) ; let \(q_{0}\) be the canonical mapping of \(N^{\prime}\) onto \(B_{0}\) , \(s_{0}\) a section of \(q_{0}\) . Set \(p_{0}(x) = s_{0}(q_{0}(x))\) for \(x \in N^{\prime}\) and extend \(p_{0}\) to T by taking \(p_{0}\) to be constant on A equal to an element of T. Then \(p = (p^{\prime}, p_{0})\) is a \(\mu\) -measurable mapping of T into the locally compact space \(B = B^{\prime} \times T\) ; it is immediate that if \(x \in N^{\prime}\) , \(y \in N^{\prime}\) , the relation \(p_{0}(x) = p_{0}(y)\) implies \(R\{x, y\}\) ; thus p meets the requirements. Moreover, if T admits a countable base, then so do each of the quotient spaces \(B_{n}\) (GT, IX, §2, No.~10, Prop. 17), therefore \(B^{\prime}\) admits a countable base, hence so does B.

PROPOSITION 3. --- Let T be a Polish locally compact space, \(\mu\) a positive measure on T, and R an equivalence relation in T. The following properties are equivalent:

\begin{enumerate}
\def\labelenumi{\alph{enumi})}
\item
  R is \(\mu\) -measurable.
\item
  There exists a sequence of mappings \(p_n: \mathrm{T} \to \mathrm{F}_n\) into Hausdorff topological spaces, such that each \(p_n\) is \(\mu\) -measurable and such that the relation \(\mathrm{R}\{x, y\}\) is equivalent to « for all \(n\) , \(p_n(x) = p_n(y)\) ».
\item
  There exists a sequence \((\mathrm{A}_{n})\) of \(\mu\) -measurable sets, saturated for R, such that for every \(x \in T\) the class of x with respect to R is the intersection of those \(A_{n}\) that contain x.
\end{enumerate}

With notations as in the statement of \(b\) , set \(p(x) = (p_n(x))\) ; the property \(b\) means that the mapping \(p\) of \(T\) into the product space \(\prod_{n} F_n\) is measurable (Ch. IV, §5, No.~3, Th. 1) and that the relation \(R\{x, y\}\) is equivalent to \(p(x) = p(y)\) ; thus \(b\) ) implies \(a\) .

Next let us show that \(c)\) implies \(b)\) . Suppose \(c)\) verified; then the characteristic functions \(\varphi_{\mathrm{A}_n}\) are \(\mu\) -measurable, and the hypothesis \(c)\) means that the relation \(\mathrm{R}\{x, y\}\) is equivalent to « for every \(n\) , \(\varphi_{\mathrm{A}_n}(x) = \varphi_{\mathrm{A}_n}(y)\) ».

Finally, let us show that \(a)\) implies \(c)\) . By Prop. 2, there exist a locally compact space B with a countable base, and a \(\mu\) -measurable mapping \(p\) of T into B, such that the relation \(\mathrm{R}\{x,y\}\) is equivalent to \(p(x)=p(y)\) . Let \((\mathrm{U}_{n})\) be a countable base for the topology of B. The sets \(\mathrm{A}_{n}=\overline{p}^{-1}(\mathrm{U}_{n})\) are \(\mu\) -measurable (Ch. IV, §5, No.~5, Prop. 7) and saturated for R; and if \(x,y\) are points of T such that \(p(x)\neq p(y)\) , there exists an index \(n\) such that \(p(x)\in\mathrm{U}_{n}\) and \(p(y)\notin\mathrm{U}_{n}\) , which means that \(x\in\mathrm{A}_{n}\) and \(y\notin\mathrm{A}_{n}\) .

\pandocbounded{\includegraphics[keepaspectratio]{images/part003_0d7fd957.jpg}}

Remark. --- If R is a \(\mu\) -measurable equivalence relation in T, the saturation for R of a compact subset of T is not necessarily \(\mu\) -measurable (Exer. 5).

THEOREM 3. --- Let T be a locally compact space with a countable base, \(\mu\) a positive measure on T, and R a \(\mu\) -measurable equivalence relation in T. Then, there exists a \(\mu\) -measurable subset S of T that intersects each class with respect to R at one and only one point (a `measurable section' for R).

We may clearly suppose that the measure \(\mu\) is bounded and that \(\mu(T) \leqslant 1\) (Ch. V, §5, No.~6, Prop. 11). We are going to define a sequence \((S_n)\) of Borel subsets (GT, IX, §6, No.~3) such that each equivalence class with respect to R intersects the union \(S'\) of the \(S_n\) in at most one point, that for every \(n\) the saturation \(T_n\) of the union of the \(S_p\) with index \(p \leqslant n\) is \(\mu\) -measurable, and that \(\mu(T - T_n) \leqslant 1/2^n\) . The saturation \(T'\) of \(S'\) will therefore be \(\mu\) -measurable and \(N = T - T'\) will have measure zero. If \(S''\) is any section of N for the relation \(R_N\) , \(S = S' \cup S''\) will meet the requirements, since \(S'\) , being a Borel set, is \(\mu\) -measurable (Ch. IV, §5, No.~4, Cor. 3 of Th. 2), and \(S''\) has measure zero.

By Prop. 2, \(\mathrm{R}\{x,y\}\) is equivalent to the relation \(p(x)=p(y)\) , where \(p\) is a \(\mu\) -measurable mapping of T into a locally compact space F. Suppose the \(\mathrm{S}_{k}\) defined for \(k\leqslant n\) . Since \(\mathrm{T}-\mathrm{T}_{n}\) is \(\mu\) -measurable and of measure \(\leqslant 1/2^{n}\) , there exists a compact subset K of \(\mathrm{T}-\mathrm{T}_{n}\) such that \(\mu\big(\mathrm{T}-(\mathrm{T}_{n}\cup\mathrm{K})\big)\leqslant 1/2^{n+1}\) and such that the restriction of \(p\) to K is continuous. Since the induced relation \(\mathrm{R}_{\mathrm{K}}\) is closed and K is metrizable, we know that there exists a Borel subset \(\mathrm{S}_{n+1}\) of K such that, in K, each point is equivalent (mod R) to one and only one point of \(\mathrm{S}_{n+1}\) (GT, IX, §6, No.~8, Th.4). Therefore \(p(\mathrm{S}_{n+1})=p(\mathrm{K})\) , which is a compact set in F; the saturation of \(\mathrm{S}_{n+1}\) for R is the inverse image \(\overline{p}^{-1}(p(\mathrm{K}))\) , which is therefore \(\mu\) -measurable (Ch. IV, §5, No.~5, Prop. 7); it is clear that this set contains K, therefore the union \(\mathrm{T}_{n+1}\) of \(\mathrm{T}_n\) and \(\overline{p}^{-1}(p(\mathrm{K}))\) is \(\mu\) -measurable, saturated for R, and is such that \(\mu(\mathrm{T} - \mathrm{T}_{n+1}) \leqslant 1/2^{n+1}\) , which completes the proof.

